
\documentclass[letter, 11pt]{article}

    %%%%%%%%%%%%%%%%%%%%%%%%%%%%%%
    \oddsidemargin  0.0in
    \evensidemargin 0.0in
    \textwidth      6.8in
    \headheight     0.0in
    \topmargin      0.0in
    \textheight=9.0in
    \parindent=0in
    %%%%%%%%%%%%%%%%%%%%%%%%%%%%%%


\renewcommand{\familydefault}{\sfdefault} % Police 

\usepackage{amsthm,amsfonts,amssymb,amsmath, xcolor, enumitem}
\usepackage{graphicx}   %To include graphics
\usepackage[small,compact]{titlesec}

\usepackage[utf8]{inputenc}

\usepackage[french]{babel}

\usepackage{fancyhdr}
\pagestyle{fancy}

\newcommand{\var}{\mathbb{V}\text{ar}} %Symboles programmés pour simplifier le texte



\usepackage{version} % Pour faire un texte troué pour les étudiants
\includeversion{etudiant}\excludeversion{professeur} %version etudiant
%\includeversion{professeur}\excludeversion{etudiants} %version prof
\newcommand{\etudiant}[1]{\begin{etudiants}#1 \end{etudiants}}
\newcommand{\prof}[1]{\begin{professeur}\color{blue} {\bf Solution :} #1 \end{professeur}}


\fancyhead[L]{STT-2100 \\ Régression}
\fancyhead[R]{A.-S. Charest, E. Reny-Nolin, A.-A. Vallée \\ Université Laval}


\begin{document}

\vspace*{0.01 cm}
\begin{center}
{\Large Exercices - Chapitre 1}
\end{center}



\begin{enumerate}[leftmargin=*, wide, label=\bf Exercice \arabic*:]

\item Pour chacune des situations suivantes, indiquez si la relation entre les variables est déterministe ou stochastique. 

\begin{enumerate}
  \item Temps d'attente dans un service  téléphonique en fonction de l'heure de l'appel.
%  \item Fréquence des orages ou des tornades
	\item Coût d'un emprunt selon le montant initial, la durée et le taux d'intérêt.
	\item Crues maximale d'un fleuve selon l'année.
	\item Volume d'un litre d'eau en fonction de la température (à pression constante).
	\item Position de la lune selon la date.
  \item Durée de vie d'une machine selon le fabricant.
  \item Nombre de décès sur les routes selon le mois.
%  \item Nombre de défauts de fabrication dans une manufacture
  \item Proportion de votes pour le Parti Libéral en fonction de l'âge, du sexe et de la circonscription.
	\item Vitesse d'atterrissage d'une masse de 1 kg lancée dans le vide à diverses hauteurs.
  \end{enumerate}

\prof{Les situations (b), (d), (e) et (i) sont déterministes, et les autres sont stochastiques.}

\vspace{0.1 in}
\item
Indiquez si les modèles suivants  sont linéaires ou non linéaires.

\begin{enumerate}
\item $Y_i=\beta_0+\beta_1\,x_i+\beta_2\,x_i^2+\varepsilon_i$
\item $Y_i=\beta_0+\beta_1\arcsin(x_i^2)+\varepsilon_i$
\item $Y_i=\beta_0\beta_1^{x_i}+\varepsilon_i$
\item $Y_i=\beta_0^3+\log(\beta_1)x_i+\varepsilon_i$
%\item $Y_i^{4/5}-1=\frac{4}{5}(\beta_0+\beta_1x_i+\varepsilon_i)$
\item $Y_i = \beta_1\,e^{x_{1i}}+\varepsilon_i$
\item $Y_i=\beta_0 + \beta_1\,x_{1i}+ \beta_2\,\text{ln}(x_{2i})+ \beta_3\,x_{3i}^2+ \varepsilon_i$
\end{enumerate}

\prof{ En statistique, un modèle est dit {\it linéaire} si l'équation peut être écrite comme une combinaison linéaire des paramètres $\beta_j$.
\begin{enumerate}
\item linéaire. 
\item linéaire. %non linéaire en $x$, mais pouvant être rendue linéaire en choisissant $z_i = \arcsin(x_i^2)$
\item non linéaire.
\item non linéaire. On pourrait se demander s'il est vraiment nécessaire d'utiliser des fonctions des paramètres ici. L'équation serait plus simple en écrivant $Y_i=\beta_0^*+\beta_1^*\,x_i+\varepsilon_i$.
\item linéaire.
\item linéaire.
\end{enumerate}}

\end{enumerate}

\end{document}

